\begin{oving}
\begin{itemize}[leftmargin=5.6em,labelwidth=5em,labelsep=0.6em,align=left]
\item[\lett] \href{https://github.com/gustavbv/Matte4C---eksamen/blob/main/oppgavesett/F08-fouriertransform/2-Bergmann-eksamener/010-TMA4125-30-35-2022-kont-oppg7-oppgave.pdf}{Matte 4N kont.\ 2022, oppg.~7} {\color{graa}(\href{https://github.com/gustavbv/Matte4C---eksamen/blob/main/oppgavesett/F08-fouriertransform/2-Bergmann-eksamener/010-TMA4125-30-35-2022-kont-oppg7-svar.pdf}{LF})}: konvolusjon av gaussfunksjoner med tabellen
\item[\lett] \href{https://github.com/gustavbv/Matte4C---eksamen/blob/main/oppgavesett/F08-fouriertransform/3-hele-oppgaven-i-pensum/046-TMA4135-2015-host-oppg2-oppgave.pdf}{TMA4135 høst 2015, oppg.~2} {\color{graa}(\href{https://github.com/gustavbv/Matte4C---eksamen/blob/main/oppgavesett/F08-fouriertransform/3-hele-oppgaven-i-pensum/046-TMA4135-2015-host-oppg2-svar.pdf}{LF})}: $f_a*f_b$ for gaussfunksjoner
\item[\middels] \href{https://github.com/gustavbv/Matte4C---eksamen/blob/main/oppgavesett/F08-fouriertransform/2-Bergmann-eksamener/009-TMA4125-2022-var-oppg6-oppgave.pdf}{Matte 4N vår 2022, oppg.~6} {\color{graa}(\href{https://github.com/gustavbv/Matte4C---eksamen/blob/main/oppgavesett/F08-fouriertransform/2-Bergmann-eksamener/009-TMA4125-2022-var-oppg6-svar.pdf}{LF})}: FT av en avkortet cosinus
\item[\middels] \href{https://github.com/gustavbv/Matte4C---eksamen/blob/main/oppgavesett/F08-fouriertransform/1-TMA4432-ovinger-og-proveeksamen/006-TMA4432-oving-8-oppg2-oppgave.pdf}{Øving 8, oppg.~2} {\color{graa}(\href{https://github.com/gustavbv/Matte4C---eksamen/blob/main/oppgavesett/F08-fouriertransform/1-TMA4432-ovinger-og-proveeksamen/006-TMA4432-oving-8-oppg2-svar.pdf}{LF})}: $f*f$ fra definisjonen, og FT av $f$ og $g$
\item[\middels] \href{https://github.com/gustavbv/Matte4C---eksamen/blob/main/oppgavesett/F08-fouriertransform/3-hele-oppgaven-i-pensum/013-TMA4130-35-2025-kont-2-sett-oppg4-oppgave.pdf}{TMA4130 kont.\ 2025 (2.\ sett), oppg.~4} {\color{graa}(\href{https://github.com/gustavbv/Matte4C---eksamen/blob/main/oppgavesett/F08-fouriertransform/3-hele-oppgaven-i-pensum/013-TMA4130-35-2025-kont-2-sett-oppg4-svar.pdf}{LF})}: FT av stykkevis funksjon og derivasjonsregelen
\item[\vanskelig] \href{https://github.com/gustavbv/Matte4C---eksamen/blob/main/oppgavesett/F08-fouriertransform/1-TMA4432-ovinger-og-proveeksamen/005-TMA4432-oving-8-oppg1-oppgave.pdf}{Øving 8, oppg.~1} {\color{graa}(\href{https://github.com/gustavbv/Matte4C---eksamen/blob/main/oppgavesett/F08-fouriertransform/1-TMA4432-ovinger-og-proveeksamen/005-TMA4432-oving-8-oppg1-svar.pdf}{LF})}: FT av $x^2$ på $\abs x<1$ (mye delvis integrasjon)
\item[\vanskelig] \href{https://github.com/gustavbv/Matte4C---eksamen/blob/main/oppgavesett/F08-fouriertransform/3-hele-oppgaven-i-pensum/012-TMA4130-35-2025-kont-oppg7-oppgave.pdf}{TMA4130 kont.\ 2025, oppg.~7} {\color{graa}(\href{https://github.com/gustavbv/Matte4C---eksamen/blob/main/oppgavesett/F08-fouriertransform/3-hele-oppgaven-i-pensum/012-TMA4130-35-2025-kont-oppg7-svar.pdf}{LF})}: skaleringsregel og produkt som konvolusjon
\item[\vanskelig] \href{https://github.com/gustavbv/Matte4C---eksamen/blob/main/oppgavesett/F08-fouriertransform/3-hele-oppgaven-i-pensum/059-TMA4135-2011-kont-oppg2-oppgave.pdf}{TMA4135 kont.\ 2011, oppg.~2} {\color{graa}(\href{https://github.com/gustavbv/Matte4C---eksamen/blob/main/oppgavesett/F08-fouriertransform/3-hele-oppgaven-i-pensum/059-TMA4135-2011-kont-oppg2-svar.pdf}{LF})}: delbrøk, FT og et vanskelig integral
\end{itemize}
\end{oving}
